% Method section. Owner: text-quality thread (claude-13); additive edits welcome.
\subsection{Setting}
A hand-written network is an ordinary feed-forward or recurrent network
$f_\theta$ whose parameter tensors are allocated at a fixed budget and
initialised to zero, after which individual entries are set by the author.
We require that (i) every layer is a standard operation (affine map, ReLU,
convolution, ReLU recurrence, mean/max pooling), so that $\theta$ loads into
any deep-learning framework unchanged, and (ii) every non-zero entry carries a
comment stating why it has its value. We report two numbers for each network:
the allocated parameter count $|\theta|$ and the number of non-zero entries
$\|\theta\|_0$.

\subsection{Input representation}
The weights are only as readable as the features they read. We therefore use
inputs whose meaning is fixed a priori and make the first embedding the
identity. For text, each character is mapped to a binary attribute vector
(Unicode general category, case, script block, punctuation class, combining
mark, a same-as-previous-character bit, and a few task-specific markers such
as the lead and tail bytes of UTF-8 mojibake). No vocabulary is involved, so
one set of weights applies to every language and script.

\subsection{Constructions}
Six constructions account for every non-zero weight in the networks reported
here.

\paragraph{Pattern kernels.} A convolution channel with weight $1$ on each
required (offset, attribute) pair and bias $-(n-1)$ fires under ReLU only when
all $n$ attributes are present: $\mathrm{relu}(\text{terminal}_{t-2} +
\text{space}_{t-1} + \text{upper}_t - 2)$ detects a sentence start in cased
scripts. A large negative weight on an attribute implements exclusion.

\paragraph{First-of-run counting.} $\mathrm{relu}(x_t - \text{same\_prev}_t)$
is $1$ only at the first character of a run of $x$, so that a separator line
of eighty dashes contributes one count rather than eighty.

\paragraph{Run-length counters.} A ReLU recurrence
$h_t=\mathrm{relu}(h_{t-1}+u_t)$ with unit self-weight and input
$u_t = \text{inc}_t - B\cdot\text{reset}_t$ ($B$ large) counts increments since
the last reset: characters since the last space (word length), since the last
sentence end (sentence length), words since the last sentence end. Because the
self-weight is exactly $1$, the recurrence has the closed form
$h_t = S_t - \min(0,\min_{k\le t} S_k)$ with $S$ the cumulative sum of $u$
(the Lindley recursion), so inference is vectorised; the step-by-step
recurrence is used for the PyTorch export and the two are tested to agree.

\paragraph{Capped detectors.} A penalty that rises from $0$ at threshold $a$
to its full value at $b$ and saturates is $\min(\mathrm{relu}(x),1) =
\mathrm{relu}(x) - \mathrm{relu}(x-1)$ with $x$ an affine function of the
pooled features scaled so that $x=1$ is ``fully bad''. The next layer multiplies
the pair by a point budget. Every penalty and reward term is one such unit, and
the sum of the named terms is the output; this is what makes
\texttt{explain()} exact rather than approximate.

\paragraph{Gates.} A detector that is meaningless for some input mode (word
length when the script has no spaces; capitalisation when the script has no
case) receives a weight of $-100$ on the rate of that mode. A soft gate
$\min(g,p)$, used to switch rate-based penalties on only once the document is
long enough for rates to be meaningful, costs one extra layer:
$p - \mathrm{relu}(p-g)$.

\paragraph{Count-then-threshold.} ``More than one script present'' is one
capped unit per script, summed with bias $-1$ and passed through ReLU; the
output is the number of extra scripts, itself capped.

\subsection{Budgets and naming}
The output is a base value plus named reward terms minus named penalty terms,
grouped into sub-scores and clipped to the target range by a final
$\mathrm{relu}(z)-\mathrm{relu}(z-100)$. Budgets are chosen in output units:
in the text-quality network a markup rate above 10\% can cost at most 20 of 100
points, mojibake at most 35, and the rewards for the shape of running prose sum
to 28, so that clean prose reaches 100 and a table reaches 0.

\subsection{Development protocol}
We develop against three sources in order: (1) hand-written samples across
scripts and junk types, each a test case with a stated reason; (2) a labelled
external set the network has never seen; (3) the worst mistakes in both
directions on (2), fixed at the level of the mechanism rather than the sample.
Every fix is recorded as a comment in the code and a line in
Section~\ref{sec:tq-devlog}. The network never sees a training signal; the
author does.
